% !TEX program = lualatex
% Transparencias de Fundamentos de las Matemáticas I
% Clase 02: Lenguaje matemático II

\documentclass[aspectratio=169,11pt]{beamer}

\usepackage{fontspec}
\usepackage[spanish,es-nodecimaldot]{babel}
\usepackage{amsmath,amssymb,mathtools}
\usepackage{booktabs}
\usepackage{csquotes}
\usepackage{xcolor}

\usetheme{Madrid}
\usecolortheme{default}
\setbeamertemplate{navigation symbols}{}
\setbeamertemplate{footline}[frame number]

\definecolor{FdMBlue}{RGB}{20,55,100}
\definecolor{FdMRed}{RGB}{130,30,30}
\definecolor{FdMGreen}{RGB}{25,100,60}

\setbeamercolor{structure}{fg=FdMBlue}
\setbeamercolor{title}{fg=white,bg=FdMBlue}
\setbeamercolor{frametitle}{fg=white,bg=FdMBlue}
\setbeamercolor{block title}{fg=white,bg=FdMBlue}
\setbeamercolor{block body}{bg=blue!3}
\setbeamercolor{alerted text}{fg=FdMRed}

\setmainfont{Latin Modern Roman}
\setsansfont{Latin Modern Sans}
\setmonofont{Latin Modern Mono}

\newcommand{\N}{\mathbb{N}}
\newcommand{\Z}{\mathbb{Z}}
\newcommand{\Q}{\mathbb{Q}}
\newcommand{\R}{\mathbb{R}}
\newcommand{\C}{\mathbb{C}}
\newcommand{\Pow}{\mathcal{P}}
\newcommand{\id}{\operatorname{id}}
\newcommand{\mcd}{\operatorname{mcd}}
\newcommand{\mcm}{\operatorname{mcm}}

\title[FdM I -- Clase 02]{Lenguaje matemático II}
\subtitle{Definiciones, enunciados y demostraciones}
\author{Antonio Falcó Montesinos}
\institute{Universidad CEU Cardenal Herrera}
\date{Fundamentos de las Matemáticas I}

\begin{document}

\begin{frame}
  \titlepage
\end{frame}

\begin{frame}{Objetivos de la clase}
\begin{itemize}
  \item Distinguir definición, enunciado, ejemplo, contraejemplo y demostración.
  \item Evaluar si una definición fija un significado preciso.
  \item Reconocer por qué muchos ejemplos no prueban un enunciado universal.
  \item Escribir una demostración elemental con hipótesis, desarrollo y cierre.
\end{itemize}
\end{frame}

\begin{frame}{Mapa de la clase}
\tableofcontents
\end{frame}

\section{Definir}

\begin{frame}{Definir}
\begin{block}{Qué hace una definición}
Una definición asigna significado preciso a un término, símbolo, objeto,
propiedad o relación dentro de un contexto determinado.
\end{block}

\begin{itemize}
  \item No se demuestra como un teorema.
  \item Debe indicar el conjunto o contexto de trabajo.
  \item Debe permitir usar el concepto sin ambigüedad.
\end{itemize}
\end{frame}

\begin{frame}{Definir}
\begin{block}{Ejemplo: número par}
Sea \(n\) un número entero. Diremos que \(n\) es par si existe un número entero
\(k\) tal que
\[
  n=2k.
\]
\end{block}

\begin{itemize}
  \item Contexto: \(n\in\Z\).
  \item Símbolo auxiliar: \(k\in\Z\).
  \item Condición: existencia de una representación \(n=2k\).
\end{itemize}
\end{frame}

\begin{frame}{Definir}
\begin{block}{Definición de subconjunto}
Sean \(A\) y \(B\) dos conjuntos. Diremos que \(A\) es un subconjunto de \(B\),
y escribiremos \(A\subseteq B\), si todo elemento de \(A\) es también elemento
de \(B\).
\end{block}

\begin{alertblock}{Control}
La notación \(A\subseteq B\) aparece después de declarar qué son \(A\) y \(B\).
\end{alertblock}
\end{frame}

\begin{frame}{Definir}
\begin{block}{Definiciones defectuosas}
\begin{itemize}
  \item Circular: usa el término que pretende definir.
  \item Incompleta: no fija el conjunto de trabajo.
  \item Ambigua: no permite decidir casos básicos.
\end{itemize}
\end{block}

\begin{exampleblock}{Ejemplo defectuoso}
``Un conjunto finito es un conjunto que tiene finitos elementos''.
\end{exampleblock}
\end{frame}

\section{Afirmar}

\begin{frame}{Afirmar}
\begin{block}{Enunciado matemático}
Un enunciado matemático es una afirmación formulada con precisión que, dentro de
un contexto dado, tiene valor de verdad.
\end{block}

\begin{itemize}
  \item ``Calcula \(2+3\)'' no es un enunciado: es una orden.
  \item ``\(x+1\)'' no es un enunciado: es una expresión.
  \item ``Para todo \(n\in\N\), \(n+1>n\)'' sí es un enunciado.
\end{itemize}
\end{frame}

\begin{frame}{Afirmar}
\begin{block}{Formas frecuentes}
\begin{itemize}
  \item Universal: para todo objeto se cumple una propiedad.
  \item Existencial: existe al menos un objeto con una propiedad.
  \item Implicación: si se cumple la hipótesis, entonces se cumple la conclusión.
  \item Equivalencia: una propiedad se cumple si y solo si se cumple otra.
\end{itemize}
\end{block}
\end{frame}

\begin{frame}{Afirmar}
\begin{block}{Clasificación rápida}
\begin{center}
\begin{tabular}{lll}
\toprule
Frase & Tipo & Comentario \\
\midrule
Sea \(A\) un conjunto & declaración & fija contexto \\
\(2\) es par & enunciado & verdadero \\
Calcula \(7+5\) & orden & no tiene valor de verdad \\
\(x+1\) & expresión & no afirma nada \\
\bottomrule
\end{tabular}
\end{center}
\end{block}
\end{frame}

\section{Ejemplos}

\begin{frame}{Ejemplos y contraejemplos}
\begin{block}{Ejemplo}
Un ejemplo es un objeto concreto que satisface una definición o propiedad.
\end{block}

\begin{block}{Contraejemplo}
Un contraejemplo a un enunciado universal es un objeto que cumple las hipótesis
pero no cumple la conclusión.
\end{block}
\end{frame}

\begin{frame}{Ejemplos y contraejemplos}
\begin{block}{Ejemplo de número par}
El número \(8\) es par porque existe \(k\in\Z\), concretamente \(k=4\), tal que
\[
  8=2k.
\]
\end{block}

\begin{block}{Contraejemplo}
La afirmación ``todo entero positivo es par'' es falsa: \(3\) es entero positivo
y no es par.
\end{block}
\end{frame}

\begin{frame}{Ejemplos y contraejemplos}
\begin{alertblock}{Regla esencial}
Un solo contraejemplo basta para refutar un enunciado universal.
\end{alertblock}

\begin{block}{Pero}
Muchos ejemplos no bastan para demostrar un enunciado universal.
\end{block}

\begin{exampleblock}{Pregunta}
¿Por qué comprobar \(2+4\), \(4+6\) y \(10+12\) no demuestra que la suma de dos
números pares cualesquiera sea par?
\end{exampleblock}
\end{frame}

\section{Demostrar}

\begin{frame}{Demostrar}
\begin{block}{Calcular, justificar, argumentar, demostrar}
\begin{itemize}
  \item Calcular: transformar expresiones para obtener un resultado.
  \item Justificar: explicar por qué un paso es válido.
  \item Argumentar: encadenar razones hacia una conclusión.
  \item Demostrar: establecer una afirmación usando definiciones, resultados previos y reglas lógicas.
\end{itemize}
\end{block}
\end{frame}

\begin{frame}{Demostrar}
\begin{block}{Estructura mínima de una demostración}
\begin{enumerate}
  \item Declarar los objetos y el contexto.
  \item Escribir las hipótesis.
  \item Indicar qué se quiere probar.
  \item Usar las definiciones pertinentes.
  \item Encadenar pasos justificados.
  \item Cerrar con la conclusión exacta.
\end{enumerate}
\end{block}
\end{frame}

\begin{frame}{Demostrar}
\begin{block}{Demostración modelo}
Demostremos que la suma de dos números pares es par.
\end{block}

Sean \(a,b\in\Z\) dos números pares. Por definición, existen \(r,s\in\Z\) tales que
\[
  a=2r
  \quad\text{y}\quad
  b=2s.
\]
Entonces
\[
  a+b=2r+2s=2(r+s).
\]
\end{frame}

\begin{frame}{Demostrar}
\begin{block}{Cierre de la demostración}
Como \(r+s\in\Z\), existe un entero \(t=r+s\) tal que
\[
  a+b=2t.
\]
Por definición de número par, \(a+b\) es par.
\end{block}

\begin{alertblock}{Observación}
El cálculo aparece dentro de la demostración, pero no sustituye al argumento.
\end{alertblock}
\end{frame}

\section{Trabajo de aula}

\begin{frame}{Trabajo de aula}
\begin{block}{Actividad 1: mejorar una frase}
Reescribe de forma rigurosa:
\begin{center}
  \Large ``Si un número acaba en cero, entonces es divisible''
\end{center}
\end{block}

\begin{itemize}
  \item ¿Qué tipo de número?
  \item ¿Qué significa ``acaba en cero''?
  \item ¿Divisible por qué número?
\end{itemize}
\end{frame}

\begin{frame}{Trabajo de aula}
\begin{block}{Una posible formulación}
Sea \(n\) un número entero escrito en base decimal. Si la última cifra decimal
de \(n\) es \(0\), entonces \(n\) es divisible por \(10\).
\end{block}

\begin{block}{Control de calidad}
\begin{itemize}
  \item El contexto está fijado.
  \item La hipótesis es reconocible.
  \item La conclusión dice exactamente qué divisibilidad se afirma.
\end{itemize}
\end{block}
\end{frame}

\begin{frame}{Trabajo de aula}
\begin{block}{Actividad 2: contraejemplos}
Da un contraejemplo para cada afirmación falsa:
\begin{enumerate}
  \item Todo número entero es positivo.
  \item Todo número impar es primo.
  \item Si \(a^2=b^2\), entonces \(a=b\), con \(a,b\in\Z\).
\end{enumerate}
\end{block}
\end{frame}

\begin{frame}{Cierre}
\begin{itemize}
  \item Una definición introduce significado; un teorema requiere prueba.
  \item Un enunciado debe tener valor de verdad en un contexto fijado.
  \item Un contraejemplo refuta una afirmación universal.
  \item Una demostración debe declarar hipótesis, usar definiciones y cerrar la conclusión.
\end{itemize}
\end{frame}

\begin{frame}{Trabajo recomendado}
\begin{itemize}
  \item Leer las secciones 1.5 a 1.10 del manual.
  \item Escribir una demostración de que la suma de dos impares es par.
  \item Revisar si todos los símbolos usados aparecen definidos antes.
\end{itemize}
\end{frame}

\end{document}
